\section{Use the quaternionic commutant in the torsion field}

This refinement uses the same inherited geometric, good-reduction, scalar
twist, compatible-representation, conductor, and height ingredients as
the pinned Version 10 source at commit
\texttt{036cebe4f1a4aa7ca89820a9b48f3629b4c364d6}. It does not replace
those ingredients by a computer check. The finite exceptional set and the
specialized local logarithm estimates are those proved in
\texttt{arithmetic\_round4.tex}. That note's anchor-nineteen construction
remains a separate fallback. This note justifies the smaller coefficient
selected for the final construction.

\subsection{An exact degree divisibility bound}

The actual construction in \texttt{manuscripts/quartic\_v10.tex},
lines 6861--7006, equips the fiber abelian surface $A/L$ with the action of
a maximal order $\mathcal O$ in the quaternion algebra of discriminant
$210$. It gives $[L:\mathbb Q]\mid576$, and every geometric
endomorphism of $A$ is defined over $L$. The non-CM argument in lines
7090--7151 takes $E=L(A[11])$ and proves $E/\mathbb Q$ Galois.
Retain that proof of normality and all of its hypotheses.

Since $11\nmid210$, the maximal order satisfies
\[
 \mathcal O/11\mathcal O\simeq M_2(\mathbb F_{11}).
\]
Its action on $V=A[11](\overline{\mathbb Q})$ is unital. Thus $V$ is
a four-dimensional module over this matrix algebra. For completeness,
its matrix units give $V=e_{11}V\oplus e_{22}V$, and $e_{12},e_{21}$
identify these summands. Both have dimension two. Every endomorphism
commuting with the matrix algebra is uniquely determined by its action
on $e_{11}V$, and every linear endomorphism of that summand extends by
the matrix units. Therefore
\[
 \operatorname{Aut}_{M_2(\mathbb F_{11})}(V)
       \simeq\operatorname{GL}_2(\mathbb F_{11}).
\]
Equivalently, $V$ is the direct sum of two standard simple modules.
No classification of an unknown Galois image is being invoked.

Since all $\mathcal O$-endomorphisms are $L$-defined, the faithful
action of $\operatorname{Gal}(E/L)$ on $V$ commutes with them. Lagrange's
theorem and the degree tower give the divisibilities
\[
 [E:L]\mid|\operatorname{GL}_2(\mathbb F_{11})|
       =(121-1)(121-11)=13200,
\]
\[
 [E:\mathbb Q]\mid576\cdot13200
       =2^{10}3^3 5^2 11.
\]
It is the divisibility, rather than only an upper bound on the numerical
degree, that controls the prime valuations of local ramification indices.
At $p=2,3,5,7$ it gives the respective bounds $(10,3,2,0)$ for
$v_p(e_E)$.

\subsection{The same local conductor proof with smaller constants}

Use exactly the fields $F=E\mathbb Q(\chi)$ and $F'$ of the source's
conductor proof, lines 7358--7404. At odd primes both added extensions
are tame. At $2$, each adds ramification index at most two. Hence
\[
 (v_2(e_{F'}),v_3(e_{F'}),v_5(e_{F'}),v_7(e_{F'}))
       \le(12,3,2,0).
\]
The source's different estimate and upper-numbering integral bound give
the conservative last-break bound $2(k+1)$ when $v_p(e_{F'})\le k$.
For the two-dimensional representation this gives Swan conductor at most
$4(k+1)$ and tame conductor at most two. At $7$ the extension is tame,
so the Swan conductor is zero. We retain these conservative bounds:
\[
 a_2\le54,\qquad a_3\le18,\qquad a_5\le14,\qquad a_7\le2.
\]
The source's companion-at-three argument still supplies the conductor
at $5$; no coefficient-characteristic conductor is omitted. At every
$p>7$ the previous bound $a_p\le2$ applies, and the conductor is
unramified outside the same set $T$.

Put
\[
 C_{\rm QM}=2^{54}3^{18}5^{14}7^2
  =2087269329221403122073600000000000000<2^{121}.
\]
The above bounds imply
\[
 N\le C_{\rm QM}R^2,\qquad \operatorname{rad}(N)\mid R.
\]
This includes a harmless extra factor of $p^2$ at the four fixed bad
primes; no subtraction of their tame contribution is needed. All other
inputs to the exponent-nineteen height proof remain unchanged.

\subsection{The height constant and anchor sixteen}

The actual proof of the exponent-nineteen estimate in lines 8915--9073
is uniform in a conductor constant $C\ge1$. Its final coefficient is
$2^{109}C^6$. For
\[
 \mathcal S=\{p:p\le41\}\cup\{109,439,2399\},\qquad
 B_{\mathcal S}=34926271784737221178290<2^{75},
\]
the argument of \texttt{arithmetic\_round4.tex} therefore gives
\[
 \mathcal H_{\rm QM}
 =2\cdot2^{109}C_{\rm QM}^6B_{\mathcal S}^{19}+128,
 \qquad h(t)\le\mathcal H_{\rm QM}M_{\mathcal S}(t)^{19}.
\]
The exact constant has bit length $2257$. A convenient elementary bound
is
\[
 \mathcal H_{\rm QM}<2^{2262},\qquad
 2262=2+109+6\cdot121+19\cdot75.
\]
This also dominates the source's CM bound $854R$, so the separate CM
branch is still covered. No change is made to the height exponent, the
character-support construction, the group theorem, or the prime-pattern
statement.

Choose
\[
 C_{\rm a}=3^{2^{15}},\qquad S_0=C_{\rm a}^2=3^{2^{16}}.
\]
The unchanged local norm-anchor and quaternion argument then gives
\[
 0\ne a-1\in2\cdot3^{2^{14}-1}\mathbb Z.
\]
Since $\log_2 3>3/2$,
\[
 \log_2|a-1|\ge1+(2^{14}-1)\log_2 3
 >24575.5>22690=10(7+2262)
 >\log_2(80\mathcal H_{\rm QM})^{10}.
\]
The order-twenty-four scalar argument consequently supplies
$h(\beta)<|a|/40$. The same Lucas comparisons at order forty and the
input comparison at order twelve reach their previous contradictions.
All completeness arguments continue to choose contact primes greater
than $76176$ and retain their previous finite exclusions.

The fixed-coefficient core count is unaffected by this coefficient
reduction. The standard decimal lift starts at $3^{2^{10}}$ and ends
at $3^{2^{16}}$, so has seven coordinates. An intermediate decimal
lift starts at $3^{2^{15}}$ and ends at $3^{2^{16}}$, so has two
coordinates. The former start exponent eighteen is not used for this
smaller endpoint. Direct decimal specialization of the fixed coefficient
is also possible without adding coordinates; it is a coefficient-writing
operation, not another mathematical reduction.

Every specialization and lift must bind its exponent to the declared IR
metadata. The parameterized tools and the adversarial parameter-binding
check distinguish exponents sixteen, nineteen and twenty-six, including
their sign certificates and exact modular coefficient contents.
